% Fichier complété par tools/complete_missing_corrections.py.
% Source : TEC/TEC-05/50_BancBalafre/50_BancBalafre.tex
% Statut : rédigé (9 question(s)).
\begin{corrigebox}[Corrigé rédigé]
\CorrigeQuestion{1}
$J_\Sigma$ est la somme des inerties de toutes les pièces ramenées à
                l'axe : $J_\Sigma=\sum J_i(\omega_i/\Omega)^2+
                \sum m_j(v_j/\Omega)^2$. L'application numérique se fait avec les rapports
                du tableau de données.
\CorrigeQuestion{2}
$E_c(\Sigma/0)=\tfrac12J_\Sigma\Omega^2$.
\CorrigeQuestion{3}
$\mathcal P_{ext}=C_m\Omega-C_{res}\Omega$ (ajouter le poids ou les
                efforts auxiliaires s'ils ont une puissance non nulle).
\CorrigeQuestion{4}
Avec les rendements $\eta_r$ et $\eta_b$, la puissance transmise est
                $\eta_r\eta_bC_m\Omega$; les pertes valent
                $P_{pertes}=(1-\eta_r\eta_b)C_m\Omega$ ou, côté résistant, la différence
                entre puissances d'entrée et de sortie.
\CorrigeQuestion{5}
$J_\Sigma\Omega\dot\Omega=
                \eta_r\eta_bC_m\Omega-C_{res}\Omega$, d'où
                $\boxed{\dot\Omega=(\eta_r\eta_bC_m-C_{res})/J_\Sigma}$.
\CorrigeQuestion{6}
Les couples moteur/résistant et les rendements sont supposés constants
                sur la phase; $J_\Sigma$ est constant. L'équation précédente impose alors
                une accélération constante.
\CorrigeQuestion{7}
Si la vitesse $\Omega_f$ doit être atteinte en $t_{max}$,
                $\alpha_{min}=\Omega_f/t_{max}$ (ou
                $(\Omega_f-\Omega_i)/t_{max}$).
\CorrigeQuestion{8}
$C_m=(J_\Sigma\alpha_{min}+C_{res})/(\eta_r\eta_b)$.
\CorrigeQuestion{9}
On retient le scénario qui maximise $J_\Sigma\alpha+C_{res}$ et on
                applique la formule précédente; le moteur choisi doit encore satisfaire
                les limites thermiques et de courant.
\end{corrigebox}
